\documentclass[a4paper]{article}
\usepackage[pdftex]{graphicx}
\usepackage[utf8]{inputenc}
\usepackage{enumerate}
\usepackage{icomma}
\usepackage{amssymb}
\usepackage{tikz}
\usepackage{href-ul}
\hypersetup{
	colorlinks=true,
	linkcolor=blue,
	urlcolor=blue}
\usepackage{geometry}
\geometry{a4paper, top=15mm, left=15mm, right=15mm, bottom=15mm,
	headsep=10mm, footskip=12mm}

\begin{document}
	\begin{tikzpicture}(10,3)
		\draw[ultra thick](2,0) --(10,0) -- (10,1.3) --(2,1.3) -- (2,0);
		\draw[fill=black](2,0)-- (0,.6) -- (2,.6) -- (2,0);
		\node at (6,.6) {\large Solution sheet from \href{http://www.okuyakl.com}{www.okuyakl.com}};
	\end{tikzpicture}
	\vspace{0.5 cm}
	
	\noindent{\bf Task 1.}\\
	Chain rule:
	$$f'(x)= {1 \over 2} (5x^2+1)^{-{1 \over 2}} \cdot 10x$$
	
	\noindent {\bf Task 1. b)}\\
	Product and chain rule:
	$$g'(x)=ln(9-x^2) - x \cdot {2x \over 9-x^2}$$
	
	\noindent {\bf Task 2. a)}\\
	The number of the logarithm must be greater than zero:
	$$2x-3 > 0 \quad \Rightarrow \quad x > 1.5 $$
	
	$$ \mathbb{D}_f= [1.5; \infty [ $$
	We examine this continuous function for limit values:
	$$
	\renewcommand{\arraystretch}{1.5}
	\begin{array}{rcl}
		\lim\limits_{x\to 1.5}{\ln(2x-3)}&=& -\infty \\
		\lim\limits_{x\to \infty}{\ln(2x-3)}&=& +\infty \\
		\Rightarrow \mathbb{W}_f &=& \mathbb{R}
	\end{array}
	$$
	
	\noindent {\bf Task 2. b)}\\
	We first calculate the zero of the graph of f:
	$$
	\renewcommand{\arraystretch}{1.5}
	\begin{array}{rcll}
		\ln(2x-3)&=& 0 &| e^{(\dots)}\\
		2x-3 &=& e^0=1 &| +3 \\
		2x &=& 4 &| :2\\
		x_0 &=& 2
	\end{array}
	$$
	Then we derive the function:
	$$ f'(x)={2 \over 2x-3}$$
	A tangent satisfies the equation of the line
	$$ y= m\cdot x + t $$
	with
	$$
	\renewcommand{\arraystretch}{1.5}
	\begin{array}{rcccl}
		y &=& f(2) &=& 0 \\
		m &=& f'(2)&=& {2 \over 4-3}=2 \\
		y=mx+t &\Rightarrow & 0=2\cdot 2 +t & & |-4 \\
		& & t &=& -4
	\end{array}
	$$
	
	The equation of the tangent is $$y=2x-4$$
	
	\noindent {\bf Task 2. c)}\\
	The function is strictly monotonically increasing in its domain because:
	$$f'(x)>0 \quad \forall \quad \textnormal{(for all)} \quad x \in \mathbb{D}_f$$
	Reason: A fraction takes on the value zero when the numerator becomes zero. In term
	$ f'(x)={2 \over 2x-3}$, however, the counter cannot become zero, no matter what you substitute for x.
	The inverse function can be determined as follows:
	$$
	\renewcommand{\arraystretch}{1.5}
	\begin{array}{rcll}
		\ln(2x-3)&=& y &| e^{(\dots)}\\
		2x-3 &=& e^y &| +3 \\
		2x &=& e^y +3 &| :2\\
		x &=& {e^y +3 \over 2} &| x\rightleftharpoons y\\
		f^{-1}(x)&=& {e^x +3 \over 2}
	\end{array}
	$$
	The definition set of an inverse function is equal to the set of values of the original function,
	so $$\mathbb{D}_{f^{-1}}= \mathbb{R}$$
	
	\noindent {\bf Task 3. a)}\\
	We derive $f(x)$ using the product and chain rule:
	$$f'(x)=e^{1-x} + x \cdot e^{1-x} \cdot (-1) = (1-x) \cdot e^{1-x}$$
	
	\noindent {\bf Task 3. b)}\\
	To investigate the monotonic behavior, we first set $f'(x)=0$:
	
	$$
	\renewcommand{\arraystretch}{1.5}
	\begin{array}{rcll}
		(1-x) \cdot e^{1-x} &=& 0 &| :e^{1-x} \quad e^{(\dots)} \quad \textnormal{can never become zero}\\
		1-x &=& 0 &| +x \\
		1 &=& x
	\end{array}
	$$
	Around the point $x=1$ we examine the monotonicity behavior of the function using a sign table:
	$$
	\renewcommand{\arraystretch}{1.5}
	\begin{array}{r|c|c|c}
		& x<1 & x=1 & x>1 \\
		\hline
		e^{1-x} & + & + & + \\
		(1-x) & + & 0 & - \\
		\hline
		f'(x) & + & 0 & - \\
		\hline
		f(x) & \nearrow & \rightarrow & \searrow \\
		& \textnormal{sms} & H & \textnormal{smf}
	\end{array}
	$$
	
	\noindent The function is strictly monotonically increasing for $x \in ]-\infty ; 1[$. \\
	The function is strictly monotonically decreasing for $x \in ]1; \infty[$.\\
	The high point at $x=1$ has the y coordinate:
	$$y=f(1)=1 \quad \Rightarrow \quad H(1|1)$$
	
	\noindent {\bf Task 3. c)}\\
	We derive the proposed antiderivative $F$ and compare it with $f$:
	$$
	\renewcommand{\arraystretch}{1.5}
	\begin{array}{rcl}
		F'(x)&=& e^{1-x} - (1-x)\cdot e^{1-x}\\
		&=& (1-(1-x))\cdot e^{1-x} \\
		&=& x \cdot e^{1-x} = f(x)
	\end{array}
	$$
	The zero of $F$ is at:
	$$
	\renewcommand{\arraystretch}{1.5}
	\begin{array}{rcl}
		-e^{1-x}- x \cdot e^{1-x} &=& 0 \\
		e^{1-x} (-1-x) &=& 0 \qquad| : e^{1-x}\\
		(-1-x) &=& 0 \\
		x &=& -1
	\end{array}
	$$
	
	\begin{center}
		\includegraphics[width=7 cm]{../../viecher/eendcomic.pdf}
		
		Here you can go back to the \href{https://www.okuyakl.de/math/m11kdL062/ae062.pdf}{Task sheet}
	\end{center}
	
	
\end{document}